\SourcePage{352}
\SGAsection{3}{半稳定约化情形。半稳定约化定理}
先给出一个具有 SGA 3 风格的辅助结果。

\begin{SGAresult}{命题}{3.1}
设 $S$ 为\Term{scheme}{概形}，$G$ 为 $S$ 上\Term{flat}{平坦}、\Term{separated}{分离}、\Term{finite-presentation}{有限表示}的\Term{commutative-group-scheme}{交换群概形}，$H$ 为 $S$ 上\Term{quasi-separated}{拟分离}且\Term{locally-finite-type}{局部有限型}的\Term{group-scheme}{群概形}，$u:G\to H$ 为 $S$-\Term{group-scheme}{群概形}的\Term{homomorphism}{同态}，$K=\operatorname{Ker}u$。对每个 $s\in S$，以 $\bar s$ 表示 $k(s)$ 的一个\Term{algebraic-closure}{代数闭包}的\Term{spectrum}{谱}，以 $\mu(s)$（相应地 $\alpha(s)$）表示 $K_{\bar s}$ 的\Term{reductive-rank}{约化秩}（相应地\Term{abelian-rank}{阿贝尔秩}），即 $(K_{\bar s})^\circ_{\mathrm{red}}$ 的最大\Term{subtorus}{子环面}（相应地最大的\Term{abelian-variety}{阿贝尔簇}\Term{quotient}{商}）的\Term{dimension}{维数}。设 $s\in S$，$t\in S$ 为 $s$ 的一个\Term{generalization}{泛化点}，$p$ 为 $k(s)$ 的\Term{characteristic-exponent}{特征指数}。

 a) 对每个与 $p$ 互素的整数 $n>0$，${}_nK$ 在 $s$ 的某\Term{neighborhood}{邻域}上是 $S$ 上\Term{etale}{平展}、\Term{finite-type}{有限型}且\Term{separated}{分离}的，并有整除关系
\[
 \operatorname{rang}({}_nK_s)\mid\operatorname{rang}({}_nK_t). \tag{3.1.1}
\]
 b) 有
\[
 \mu(s)+2\alpha(s)\leq\mu(t)+2\alpha(t). \tag{3.1.2}
\]
 c) 若 $(K_{\bar t})^\circ_{\mathrm{red}}$ 是\Term{unipotent}{幺幂}的，则 $(K_{\bar s})^\circ_{\mathrm{red}}$ 亦然。

 d) 假设 $G_s$ 的\Term{unipotent-rank}{幺幂秩}为零，$H$ 在 $S$ 上\Term{smooth}{光滑}，且 $u_t$ 是\Term{isogeny}{同源}。则在 $s$ 的某个\Term{open-neighborhood}{开邻域} $U$ 上，$u$ 是\Term{flat}{平坦}且\Term{quasi-finite}{拟有限}的，因而 $K|U$ 也是\Term{flat}{平坦}且\Term{quasi-finite}{拟有限}的。

 e) 在 d) 的条件下，若 $u_t$ 还是\Term{monomorphism}{单态射}（从而是 $G_t$ 到 $H_t$ 的某个\Term{open-subgroup}{开子群}的\Term{isomorphism}{同构}），则 $u_s$ 是 $G_s$ 到 $H_s$ 的某个\Term{open-subgroup}{开子群}的\Term{isomorphism}{同构}；
\SourcePage{353}
而且，若 $S$ 还是\Term{irreducible}{不可约}的，且 $t$ 为其\Term{generic-point}{泛点}，则存在 $s$ 的一个\Term{open-neighborhood}{开邻域} $U$，使 $u$ 诱导 $G|U$ 到 $H|U$ 的某个\Term{open-subgroup}{开子群}的\Term{isomorphism}{同构}。
\end{SGAresult}
\begin{proof}
 a) 可以假设 $n$ 在 $S$ 上\Term{invertible}{可逆}；此时 $n\cdot\operatorname{id}_G$ 在每个\Term{fiber}{纤维}上都\Term{etale}{平展}（SGA 3 XV 1.3）。由于 $G$ 在 $S$ 上\Term{flat}{平坦}且\Term{finite-presentation}{有限表示}，由逐\Term{fiber}{纤维}的\Term{etale}{平展}判据 EGA IV 17.8.2 可知，$n\cdot\operatorname{id}_G$ 是\Term{etale}{平展}的，因而其\Term{kernel}{核} ${}_nG$ 是\Term{etale}{平展}的。又因 $G$ 在 $S$ 上\Term{separated}{分离}且\Term{finite-presentation}{有限表示}，显然 ${}_nG$ 也具有这两个性质。同样可见，${}_nH$ 在 $S$ 上\Term{locally-finite-type}{局部有限型}、\Term{quasi-separated}{拟分离}且\Term{formally-unramified-old}{形式净}\footnote{至少当 $H$ 交换时如此；在\Term{generic}{一般}情形中，${}_nH$ 沿其\Term{unit-section}{单位截面}在 $S$ 上仍是\Term{net}{净}的，这对我们的目的已经足够。}，因此它的\Term{unit-section}{单位截面}是\Term{open-immersion}{开浸入}（EGA IV 17.4.2 b')）。故 ${}_nu:{}_nG\to{}_nH$ 的\Term{kernel}{核}是 ${}_nG$ 的\Term{open-subscheme}{开子概形}，从而在 $S$ 上\Term{etale}{平展}；它在 ${}_nG$ 中是\Term{retrocompact}{相对拟紧}的，因而在 $S$ 上\Term{finite-presentation}{有限表示}，又因 ${}_nG$ \Term{separated}{分离}，它在 $S$ 上也\Term{separated}{分离}。这个\Term{kernel}{核}显然就是 ${}_nK$。
公式 (3.1.1) 由 ${}_nK$ 的上述性质得出：作\Term{base-change}{基变换}化到 $S$ 为以 $s$ 为\Term{closed-point}{闭点}的\Term{henselian}{Hensel}\Term{local-scheme}{局部概形}的情形，再利用 ${}_nK$ \Term{resolution}{分解}为 ${}_nK^{\mathrm f}$ 与一个\Term{special-fiber}{特殊纤维}为空的\Term{scheme}{概形}之和（参见 2.2.3），即可看出。事实上，这一论证更精确地表明，\Term{geometric-fiber}{几何纤维}之间存在\Term{monomorphism}{单态射}
\[
 {}_nK(\bar s)\lhook\joinrel\longrightarrow{}_nK(\bar t).
\]
\SourcePage{354}
 b) 取整数 $n>0$，使其既与 $k(s)$ 的\Term{characteristic-exponent}{特征指数}互素，又与 $K$ 在 $s,t$ 处的\Term{geometric-fiber}{几何纤维}模去其\Term{reduced-neutral-component}{约化单位分支}所得\Term{quotient}{商}的\Term{rank}{秩}互素。由\Term{commutative-algebraic-group}{交换代数群}的已知结构，(3.1.1) 两边分别等于以 (3.1.2) 的两边为指数的 $n$ 次幂。因此 b) 由 a) 得出。

 c) 假设 $(K_{\bar t})^\circ_{\mathrm{red}}$ \Term{unipotent}{幺幂}，\Term{equivalence}{等价}于 $\mu(t)=\alpha(t)=0$；对 $K_s$ 也有同样的说法。因此断言 c) \Term{form}{形式}地由 b) 得出。

 d) 假设 $u_t$ 为\Term{isogeny}{同源}，则 $K_t$ \Term{finite}{有限}，故满足 c) 的条件，从而 $(K_{\bar s})^\circ_{\mathrm{red}}$ 是\Term{unipotent}{幺幂}的。由于 $G_s$ 的\Term{unipotent-rank}{幺幂秩}为零，可知 $K_s$ \Term{finite}{有限}。另一方面，$u_t$ 为\Term{isogeny}{同源}意味着 $\dim G_t=\dim H_t$，而 $G,H$ \Term{flat}{平坦}且\Term{locally-finite-presentation}{局部有限表示}意味着 $\dim G_s=\dim G_t$、$\dim H_s=\dim H_t$（SGA 3 VI$_{\mathrm B}$ 4.3），故 $\dim G_s=\dim H_s$。由 $\dim\operatorname{Ker}u_s=0$ 可得 $\dim H_s=\dim u_s(G_s)$（SGA 3 VI$_{\mathrm B}$ 1.2）；由于 $H_s$ \Term{smooth}{光滑}，这\Term{equivalence}{等价}于 $u_s$ \Term{flat}{平坦}（SGA 3 VI$_{\mathrm A}$ 5.6）。因此存在 $s$ 的一个\Term{open-neighborhood}{开邻域} $U$，使 $u$ 在 $U$ 上的\Term{unit-section}{单位截面}各点处\Term{flat}{平坦}且\Term{quasi-finite}{拟有限}（SGA 3 VI$_{\mathrm B}$ 2.2 (i) 与 2.5 (ii)）。与 $u_s$ 的情形一样，可知对每个 $s'\in U$，$u_{s'}$ 都\Term{flat}{平坦}；于是 $u|G_U$ \Term{flat}{平坦}且\Term{quasi-finite}{拟有限}，因为它在每个\Term{fiber}{纤维}上都如此。

 e) 对第一个断言，作\Term{base-change}{基变换}即可立即化到 $S$ 是以 $t$ 为\Term{generic-point}{泛点}的\Term{integral-scheme}{整概形}的情形，因而只须证明第二个断言。后者立即由 d) 及下述\Term{lemma}{引理}得出。
\end{proof}
\SourcePage{355}
\begin{SGAresult}{\Term{lemma}{引理}}{3.1.3}
设 $S$ 为\Term{scheme}{概形}，$K$ 为 $S$ 上\Term{quasi-finite}{拟有限}、\Term{finite-presentation}{有限表示}、\Term{separated}{分离}且\Term{flat}{平坦}的 $S$-\Term{scheme}{概形}，$s\in S$，$t\in T$ 为 $s$ 的一个\Term{generalization}{泛化点}。则有
\[
 \operatorname{rang}K_s\leq\operatorname{rang}K_t.
\]
若对 $S$ 的每个\Term{maximal-point}{极大点} $m$ 都有 $\operatorname{rang}K_m\leq1$，则 $K\to S$ 是\Term{open-immersion}{开浸入}（从而若它有\Term{section}{截面}，便是\Term{isomorphism}{同构}）。
\end{SGAresult}
\begin{proof}
对第一个断言，作\Term{strict-localization}{严格局部化}，化到 $S$ 为\Term{strictly-local}{严格局部}的情形。将 $K$ \Term{resolution}{分解}为一个\Term{finite}{有限}\Term{scheme}{概形} $K^{\mathrm f}$ 与一个\Term{special-fiber}{特殊纤维}为空的\Term{scheme}{概形}之和，便有
\[
 \operatorname{rang}K_s=\operatorname{rang}K_s^{\mathrm f}
 =\operatorname{rang}K_t^{\mathrm f}\leq\operatorname{rang}K_t;
\]
其中第二个等号来自 $K^{\mathrm f}$ 在 $S$ 上\Term{flat}{平坦}、\Term{finite}{有限}且\Term{finite-presentation}{有限表示}，因而在 $S$ 上\Term{locally-free}{局部自由}。

对第二个断言，上述论证已表明，对每个 $s\in S$ 都有 $\operatorname{rang}K_s\leq1$；故 $K_s$ 为空或\Term{isomorphic}{同构}于 $\operatorname{Spec}k(s)$。因此 $K\to S$ 是\Term{monomorphism}{单态射}（SGA 3 VI$_{\mathrm B}$ 2.11），又因它\Term{flat}{平坦}且\Term{finite-presentation}{有限表示}，便是\Term{open-immersion}{开浸入}（EGA IV 17.9.1），证明完毕。
\end{proof}
\textbf{注记 3.1.4。}在 3.1 a) 及命题 3.1 的其余部分中，当 $n$ 为 $p$ 的幂时（$p>1$），没有给出任何关于 ${}_nK_s$ 的信息。至少在 $G$ \Term{smooth}{光滑}（且 $S$ \Term{locally-noetherian}{局部诺特}或 $H$ 在 $S$ 上\Term{finite-presentation}{有限表示}）时，可以用与 a) 的证明实质相同的方法给出这样的信息：
\SourcePage{356}
化到 $S$ 为\Term{complete-local-noetherian}{完备诺特局部}的情形，并使用 SGA 3 XV 1.6 c)、X 3.2、IX 3.6、IX 5.2，即得
\[
 \operatorname{rang\,mult}({}_nK_s)\mid
 \operatorname{rang\,mult}({}_nK_t). \tag{3.1.4.1}
\]
这里 $\operatorname{rang\,mult}$ 表示所考察的\Term{commutative-algebraic-group}{交换代数群}的最大\Term{multiplicative-type}{乘法型}\Term{subgroup}{子群}的\Term{rank}{秩}（SGA 3 XVII 7.2.1）；注意，\Term{multiplicative-type}{乘法型}的\Term{torsion}{挠}\Term{algebraic-group}{代数群}是\Term{finite}{有限}的。公式 (3.1.4.1) 对每个整数 $n\geq1$ 都成立。由 SGA XVII 4.6.1 (vi) 可得，在 c) 中，若 $K_t$ 是\Term{unipotent}{幺幂}的（SGA 3 VII 1.1），则 $K_s$ 亦然。

\begin{SGAresult}{命题}{3.2}
设 $S$ 为\Term{locally-noetherian}{局部诺特}、\Term{regular}{正则}、一维\Term{integral-scheme}{整概形}，$K$ 为其\Term{rational-function-field}{有理函数域}，$A_K$ 为 $K$ 上的\Term{abelian-scheme}{阿贝尔概形}，$A$ 为它在 $S$ 上的\Term{neron-model}{Néron 模型}；这是 $S$ 上\Term{smooth}{光滑}且\Term{separated}{分离}的\Term{group-scheme}{群概形}（若 $S$ \Term{noetherian}{诺特}，则还在 $S$ 上\Term{quasi-projective}{拟射影} [23，VIII 2，$3^\circ$]）。以下条件\Term{equivalence}{等价}：

 (i) 对每个 $s\in S$，$A_s^\circ$ 的\Term{unipotent-rank}{幺幂秩}为零，即 (2.1.5) 是\Term{abelian-scheme}{阿贝尔概形}被\Term{torus}{环面}的\Term{group-extension}{扩张}。

 (ii) 存在 $S$ 上\Term{smooth}{光滑}、\Term{separated}{分离}且\Term{finite-type}{有限型}的\Term{group-scheme}{群概形} $G$，它延拓 $A_K$，且其\Term{fiber}{纤维}的\Term{unipotent-rank}{幺幂秩}均为零。

此外，若 $G$ 如 (ii) 中所述，则在\Term{generic-fiber}{一般纤维}上诱导给定\Term{isomorphism}{同构}的唯一\Term{morphism}{态射} $G\to A$ 是 $G$ 到 $A$ 的某个\Term{open-subgroup}{开子群}的\Term{isomorphism}{同构}，因而诱导\Term{neutral-component}{单位分支}的\Term{isomorphism}{同构} $G^\circ\xrightarrow{\sim}A^\circ$。
\end{SGAresult}
\SourcePage{357}
\textbf{3.2.1。}先指出，$A_K$ 的\Term{neron-model}{Néron 模型} $A$ 的定义与 1.1 中回顾的\Term{local-scheme}{局部}情形相同。这个\Term{neron-model}{Néron 模型}的存在性，以及它在 $S$ 上\Term{separated}{分离}且\Term{finite-type}{有限型}这一事实，都是众所周知的，而且实际上可立即归结到\Term{local-scheme}{局部}情形。当然，对每个\Term{closed-point}{闭点} $s\in S$，$S'=\operatorname{Spec}(\mathscr O_{S,s})$ 是\Term{trait}{离散赋值环谱}；$A\times_SS'$ 正是 $A_K$ 在 $S'$ 上的\Term{neron-model}{Néron 模型}，其含义与 1.1 中所考虑的相同。
\begin{proof}[命题 3.2 的证明]
(i) 显然蕴含 (ii)，因为只须取 $G=A$。故只剩下证明：若 $G$ 如 (ii) 中所述，则 $G\to A$ 是\Term{open-immersion}{开浸入}。这正是 3.1 e) 的一个特例。
\end{proof}
\begin{SGAresult}{推论}{3.3}
沿用 3.2 的记号，假设\Term{neron-model}{Néron 模型}的\Term{fiber}{纤维}的\Term{neutral-component}{单位分支}的\Term{unipotent-rank}{幺幂秩}为零。设 $S'$ 为满足与 $S$ 相同假设的\Term{scheme}{概形}，$K'$ 为其\Term{rational-function-field}{有理函数域}，$f:S'\to S$ 为\Term{dominant-morphism}{支配态射}，$A_{K'}=A_K\otimes_KK'$ 为 $A\times_SS'$ 的\Term{generic-fiber}{一般纤维}，$A'$ 为 $A_{K'}$ 的\Term{neron-model}{Néron 模型}，并考虑在\Term{generic-fiber}{一般纤维}上诱导\Term{identity-map}{恒等映射}的唯一\Term{morphism}{态射}
\[
 u:A\times_SS'\longrightarrow A'.
\]
则 $u$ 是\Term{open-immersion}{开浸入}，因而诱导\Term{neutral-component}{单位分支}的\Term{isomorphism}{同构}
\[
 A^\circ\times_SS'\xrightarrow{\sim}A'^\circ. \tag{3.3.1}
\]
\end{SGAresult}
\begin{proof}
只须在 $S'$ 的情形应用 3.2，取 $G=A\times_SS'$ 即可。
\end{proof}
\SourcePage{358}
\textbf{3.3.2。}由此可见，3.2 的\Term{equivalence}{等价}性质蕴含\Term{neron-model}{Néron 模型}的一种部分“\Term{stable}{稳定}性”。须注意，当 $A_K$ 没有\Term{good-reduction}{好约化}，即 $A$ 不是\Term{abelian-scheme}{阿贝尔概形}时，并非对 3.3 中的每个 $S'$，\Term{morphism}{态射} $A\times_SS'\to A'$ 本身都是\Term{isomorphism}{同构}，即在\Term{fiber}{纤维}的\Term{connected-component-group}{连通分支群}上诱导\Term{isomorphism}{同构}（即便 $S'$ 在 $S$ 上\Term{finite}{有限}，且 $K'/K$ \Term{separable}{可分}，也并非如此……；关于这些\Term{connected-component-group}{连通分支群}在\Term{base-extension}{基扩张}下的行为，参见下文 11.10）。这些观察，再加上下文 3.9，说明了下述由 Mumford 引进的术语的合理性（据说它也与他在书 [17] 中研究的同名概念相协调）。

\begin{SGAresult}{定义}{3.4}
沿用 3.2 的记号。若满足 3.2 的\Term{equivalence}{等价}条件 (i)、(ii)，则称 $K$ 上的\Term{abelian-scheme}{阿贝尔概形} $A_K$ 在 $S$ 上有\Term{semistable-reduction}{半稳定约化}。
\end{SGAresult}
\textbf{3.4.0。}注意，采用 (ii) 的\Term{form}{形式}时，这个条件并不涉及\Term{neron-model}{Néron 模型}的概念。另一方面，例如由 3.2.1 可知，作\Term{localization}{局部化}便立即归结到 $S$ 为\Term{trait}{离散赋值环谱}的情形。而且，\Term{neron-model}{Néron 模型}的构造与从 $S$ 到其\Term{henselization}{Hensel 化}或\Term{strict-henselization}{严格 Hensel 化}的\Term{base-change}{基变换}可交换，因此这个概念甚至可归结到 $S$ 为\Term{strictly-local}{严格局部}\Term{trait}{离散赋值环谱}的情形。
\SourcePage{359}
\begin{SGAresult}{命题}{3.5（\Term{semistable-reduction}{半稳定约化}的 Galois 判据）}
假设 $S$ 是\Term{trait}{离散赋值环谱}。\Term{abelian-scheme}{阿贝尔概形} $A_K$ 的以下条件\Term{equivalence}{等价}（记号为 (2.5.3) 中的记号，应用于将 $(A_K,S)$ 换为 $S$ 的\Term{henselization}{Hensel 化}后得到的情形）：

 (i) $A_K$ 在 $S$ 上有\Term{semistable-reduction}{半稳定约化}（3.4）。

 (ii) 用 (2.5.3)、(2.5.4) 的记号，有
\[
 V\subset V^\perp.
\]
 (iii) 对\Term{tate-module}{Tate 模}的\Term{toric-part}{环面部分} $W$，有
\[
 W=V^\perp.
\]
 (iv) \Term{inertia-group}{惯性群} $I$ 在 $U=T_\ell(A(\overline K))$ 上的\Term{action}{作用}是二阶\Term{unipotent}{幺幂}的（即存在 $U$ 的一个在 $I$ 下\Term{stable}{稳定}的\Term{submodule}{子模} $U'$，使 $I$ 在 $U'$ 和 $U/U'$ 上均\Term{trivial-action}{平凡作用}）。
\end{SGAresult}
\begin{proof}
可以假设 $S$ 为\Term{henselian}{Hensel}的（3.4.0）。用 2.1.9 的记号，(i) 意味着
\[
 \lambda=0;
\]
由 (2.5.5)，这\Term{equivalence}{等价}于 $V=V\cap V^\perp$，即 (ii)。由\Term{orthogonality-theorem}{正交性定理} 2.4，$W=V\cap V^\perp$，故 (ii) $\Longleftrightarrow$ (iii)。此外，(iv) 也意味着 $I$ 在 $U^I=V$ 上\Term{trivial-action}{平凡作用}；由\Term{duality}{对偶性}，这意味着它在 $V^\perp$ 上\Term{trivial-action}{平凡作用}，即 $V^\perp\subset U^I=V$，这就是 (ii)。证明完毕。
\end{proof}
\textbf{注记 3.5.1。}条件 (ii)、(iv) 显然只依赖于 $I$ 在 $\mathbb Q_\ell$-\Term{vector-space}{向量空间} $U\otimes_{\mathbb Z_\ell}\mathbb Q_\ell$ 上的\Term{action}{作用}，本来也可以用这个\Term{action}{作用}来表述。还应指出，这些条件显然在任意\Term{trait}{离散赋值环谱}的\Term{morphism}{态射}
\SourcePage{360}
$S'\to S$ 所给出的\Term{base-change}{基变换}下保持成立。最后，若 $A_K$ 满足 (iv)，则它的每个\Term{abelian-subscheme}{阿贝尔子概形}、每个\Term{abelian-quotient-scheme}{阿贝尔商概形}也显然满足 (iv)；特别地，每个与 $A$ \Term{isogenous}{同源}的\Term{abelian-scheme}{阿贝尔概形}（例如 $A'_K$）都满足它。

\begin{SGAresult}{推论}{3.5.2}
设 $g$ 为 $\mathbb Z_\ell(1)=T_\ell(\overline k^{\,*})$ 的一个\Term{topological-generator}{拓扑生成元}。3.5 的\Term{equivalence}{等价}条件还\Term{equivalence}{等价}于下述条件：

 (v) \Term{inertia-group}{惯性群} $I$ 在 $U=T_\ell(A(\overline K))$ 上的\Term{action}{作用}通过 $I$ 的最大\Term{pro-l-finite}{pro-$\ell$-有限}\Term{quotient}{商} $I(\ell)$ \Term{resolution}{分解}；这个\Term{quotient}{商}\Term{canonically-isomorphic}{典范同构}于 $\mathbb Z_\ell(1)$（第 I 讲\SGAsourcegap{原 PDF 第360页此处\Term{lie-bracket}{括号}中仅印有“I”，其后的交叉引注留空。}），且\Term{generator}{生成元} $g$ 在 $U$ 上的\Term{action}{作用} $g_U$ 满足
\[
 (1-g_U)^2=0. \tag{3.5.3}
\]
\end{SGAresult}
\begin{proof}
事实上，对于\Term{profinite-group}{profinite 群} $I$ 在一个 $\ell$ 进\Term{vector-space}{向量空间}上的任意\Term{continuous}{连续}\Term{unipotent}{幺幂}\Term{action}{作用}，立即可见，$I$ 的像是\Term{pro-l-group}{pro-$\ell$ 群}。由此即得 (v) 与 (iv) \Term{equivalence}{等价}。
\end{proof}
现在证明本讲的核心结果。
\begin{SGAresult}{\Term{theorem}{定理}}{3.6（\Term{semistable-reduction}{半稳定约化}\Term{theorem}{定理}）}
设 $S$ 为一维\Term{noetherian}{诺特}、\Term{regular}{正则}、\Term{connected}{连通}\Term{scheme}{概形}，$K$ 为 $S$ 的\Term{rational-function-field}{有理函数域}，$A_K$ 为 $K$ 上的\Term{abelian-scheme}{阿贝尔概形}。则存在 $K$ 的一个\Term{finite-galois-extension}{有限 Galois 扩张} $K'$，使\Term{abelian-scheme}{阿贝尔概形} $A_{K'}=A_K\otimes_KK'$ 在 $S$ 于 $K'$ 中的\Term{normalization}{正规化} $S'$ 上有\Term{semistable-reduction}{半稳定约化}（3.4）。
\end{SGAresult}
\SourcePage{361}
\begin{proof}
注意，$S'$ 是一维\Term{noetherian}{诺特}、\Term{regular}{正则}、\Term{connected}{连通}\Term{scheme}{概形}，所以所述结论是有意义的。为证明 3.6，选取 $K$ 的一个\Term{separable-closure}{可分闭包} $\overline K$，并将\Term{extension}{扩张} $\overline K$ 写成其\Term{finite-galois-subextension}{有限 Galois 子扩张} $K_i$ 的\Term{inductive-limit}{归纳极限}。设 $S_i$ 为 $S$ 在 $K_i$ 中的\Term{normalization}{正规化}，$U$ 为 $S$ 中使 $A_K$ 有\Term{good-reduction}{好约化}的最大\Term{open-subset}{开集}。由于 $S$ \Term{noetherian}{诺特}、\Term{connected}{连通}且为一维，且 $U\ne\varnothing$，可知 $T=S-U$ 是 $S$ 的\Term{finite}{有限}个\Term{closed-point}{闭点}所成的\Term{set}{集合}，并且对每个 $t\in T$，$\mathscr O_{S,t}$ 是\Term{discrete-valuation-ring}{离散赋值环}。只须对每个 $t\in T$ 找到一个指标 $i=i(t)$，使 $A_{K_i}$ 在 $S_i$ 中 $t$ 上方的各点处都有\Term{semistable-reduction}{半稳定约化}；因为此时取各\Term{extension}{扩张} $K_{i(t)}$（$t\in T$）的\Term{compositum}{合成域}为 $K'$ 即可。这样便化到了 $S$ 为\Term{trait}{离散赋值环谱}的情形。
另一方面，\Term{group}{群} $\operatorname{Gal}(K_i/K)$ 置换 $S_i$ 中位于 $S$ 的\Term{closed-point}{闭点} $s$ 上方的各点，所以只须找到 $K'$，使 $A_{K'}$ 在\Term{normalization}{正规化} $S'$ 的某一点处有\Term{semistable-reduction}{半稳定约化}。此外，这里并不需要 $K'$ 在 $K$ 上为 Galois 的；只须找到满足前一条件的 $K$ 的一个\Term{finite-etale-extension}{有限平展扩张} $K'$，因为这个\Term{extension}{扩张}的\Term{galois-closure}{Galois 闭包}也会满足同一条件。因此，利用判据 3.5 (iv)，可见 3.6 归结为以下断言。
\end{proof}
\begin{SGAresult}{推论}{3.7}
再次假设 $S$ 是\Term{trait}{离散赋值环谱}。则存在 $\pi$ 的一个\Term{open-subgroup}{开子群} $\pi'$，使 $I'=I\cap\pi'$ 在 $U=T_\ell(A(\overline K))$ 上的\Term{action}{作用}是二阶\Term{unipotent}{幺幂}的；换言之，存在 $I$ 的一个\Term{open-subgroup}{开子群} $I'$，其在 $U$ 上的\Term{action}{作用}是二阶\Term{unipotent}{幺幂}的。
\end{SGAresult}
\SourcePage{362}
\begin{proof}
有一个与 $\pi$ 的\Term{action}{作用}相容的\Term{canonical-isomorphism}{典范同构}
\[
 H^1(A_{\overline K},\mathbb Z_\ell)\xrightarrow{\sim}
 \operatorname{Hom}(T_\ell(A(\overline K)),\mathbb Z_\ell). \tag{*}
\]
它表明，断言 3.7 \Term{equivalence}{等价}于关于 $I$ 在 (*) 左边的\Term{action}{作用}的同一断言。这个断言已在第 III 讲中证明（使用了二维\Term{excellent}{优良}\Term{scheme}{概形}的\Term{resolution-of-singularities}{奇点解消}）。
\end{proof}
\begin{SGAresult}{推论}{3.8}
3.5 的\Term{equivalence}{等价}条件 (i)--(iv) 还\Term{equivalence}{等价}于以下条件：

 (vi) $I$ 在 $U=T_\ell(A(\overline K))$ 上的\Term{action}{作用}是\Term{unipotent}{幺幂}的。
\end{SGAresult}
\begin{proof}
事实上，为证明 (v) 蕴含 (vi)，注意如下事实：若\Term{group}{群} $I$ 在\Term{characteristic-zero}{特征零}的域（这里是 $\mathbb Q_\ell$）上的\Term{finite-dimensional-vector-space}{有限维向量空间} $E$ 上作\Term{unipotent}{幺幂}\Term{action}{作用}，且它有一个以 $m$ 阶\Term{action}{作用}的\Term{finite-index-subgroup}{有限指数子群} $I'$（这里 $m=2$），则 $I$ 本身也以 $m$ 阶\Term{action}{作用}。事实上，因为 $I'$ 在 $I$ 中的指数\Term{finite}{有限}，它所定义的\Term{algebraic-group}{代数群} $G'\subset\operatorname{GL}(E)$ 在 $I$ 所定义的\Term{algebraic-group}{代数群} $G$ 中的指数\Term{finite}{有限}。然而 $G$ 是\Term{unipotent}{幺幂}的，因而\Term{connected}{连通}（这里\Term{characteristic-zero}{特征零}！），故 $G=G'$，从而 $I$ 与 $I'$ 有相同的“\Term{unipotence-level}{幺幂阶数}”。
\end{proof}
还可得到 3.3 的一个逆命题。
\begin{SGAresult}{推论}{3.9}
记号与 3.2 相同。3.2 的条件 (i)、(ii) 还\Term{equivalence}{等价}于下述条件：

 (iii) 对 $K$ 的每个\Term{etale-extension}{平展扩张} $K'$，以 $S'$ 表示 $S$ 在 $K'$ 中的\Term{normalization}{正规化}，以 $A'$ 表示 $A_{K'}=A\otimes_SK'=A_K\otimes_KK'$ 的\Term{neron-model}{Néron 模型}，则\Term{canonical-morphism}{典范态射}
\SourcePage{363}
\[
 A\times_SS'\longrightarrow A'
\]
在\Term{neutral-component}{单位分支}上诱导\Term{isomorphism}{同构}。
\end{SGAresult}
\begin{proof}
条件的必要性实际上已在 3.3 中以更强的\Term{form}{形式}见到。为证明充分性，选取 3.6 中的 $K'$。由于 $A'^\circ$ 的\Term{fiber}{纤维}的\Term{unipotent-rank}{幺幂秩}为零，$(A\times_SS')^\circ=A^\circ\times_SS'$ 的\Term{fiber}{纤维}也如此，因而 $A^\circ$ 的\Term{fiber}{纤维}也如此，因为前者是后者经\Term{base-field-extension}{基域扩张}得到的。
\end{proof}

\SGAsection{4}{应用于 Serre--Tate 的一个猜想及导子}
\textbf{4.1。}假设 $S$ 为\Term{henselian}{Hensel}的，并沿用 (2.5.3) 的记号，其中 $\ell\ne p$\footnote{$\ell\ne p$ 的情形将在下文第5、6节研究。}。为研究 $\pi=\operatorname{Gal}(\overline K/K)$ 及其\Term{inertia-subgroup}{惯性子群} $I$ 在 $U=T_\ell(A(\overline K))$ 上的\Term{action}{作用}，我们将使用以下事实：存在 $I$ 的一个\Term{open-subgroup}{开子群} $I'$，并且可令它在 $\pi$ 中\Term{normal-subgroup}{正规}，使 $I'$ 在 $U$ 上的\Term{action}{作用}为二阶\Term{unipotent}{幺幂}，即有
\[
 U^{I'}\supset(U^{I'})^\perp.
\]
显然，$U^{I'}$ 此时不依赖于 $I'$ 的具体选择。可以称它为 $U$ 的\Term{essentially-fixed-submodule}{本质固定子模}，记为 $U^{\mathrm{ef}}$；其\Term{orthogonal}{正交补}相应地称为 $U$ 的\Term{essentially-toric-submodule}{本质环面子模}，记为 $U^{\mathrm{et}}$。这样得到三级\Term{filtration}{滤过}
\[
 U\supset U^{\mathrm{ef}}\supset U^{\mathrm{et}}\supset0, \tag{4.1.1}
\]
显然在 $\pi$ 的\Term{action}{作用}下\Term{invariant}{不变}。顺便指出，
\SourcePage{364}
$I$ 中存在一个在 $U$ 上作\Term{unipotent}{幺幂}\Term{action}{作用}（或\Term{equivalence}{等价}地，作二阶\Term{unipotent}{幺幂}\Term{action}{作用}）的最大\Term{subgroup}{子群}：它就是由作\Term{unipotent}{幺幂}\Term{action}{作用}的元素组成的\Term{subgroup}{子群} $I'$，也就是 $I$ 在 $U$ 的\Term{filtration}{滤过} (4.1.1) 的\Term{associated-graded}{相伴分次对象}上的\Term{representation}{表示}的\Term{kernel}{核}；在这里，也等于 $I$ 在 $U^{\mathrm{ef}}$ 上所诱导的\Term{representation}{表示}的\Term{kernel}{核}。由\Term{semistable-reduction}{半稳定约化}的 Galois 判据（3.5），$I$ 的这个\Term{subgroup}{子群} $I'$ 对应于 $\widetilde K$（$K$ 的最大\Term{unramified-extension}{非分歧扩张}）的最小\Term{galois-extension}{Galois 扩张} $\widetilde K'$，使 $A_{\widetilde K'}$ 有\Term{semistable-reduction}{半稳定约化}。这个描述同时表明，$I$ 的\Term{subgroup}{子群} $I'$ 不依赖于\Term{prime-number}{素数} $\ell\ne p$ 的选择。

注意，知道 $U^{\mathrm{ef}}$ 上、因而知道 $U^{\mathrm{et}}$ 上的\Term{representation}{表示}，便能在相差\Term{isogeny}{同源}的意义下知道 $\pi$ 在 $U/U^{\mathrm{ef}}$ 上的\Term{representation}{表示}；这是因为\Term{polarization-pairing}{极化配对} (2.5.1) 诱导\Term{isogeny}{同源}
\[
 U/U^{\mathrm{ef}}\longrightarrow(U^{\mathrm{et}})^\vee(1)
 \qquad\text{为同源}, \tag{4.1.2}
\]
这里 $(1)$ 表示以 $\mathbb Z_\ell(1)=T_\ell(\overline K^{\,*})$ 作\Term{tate-twist}{Tate 扭转}；上述\Term{homomorphism}{同态}显然与 $\pi$ 的\Term{action}{作用}相容。（注意，有一个\Term{canonical-isomorphism}{典范同构} $U/U^{\mathrm{ef}}\xrightarrow{\sim}(U'^{\mathrm{et}})^\vee(1)$，其中 $U'=T_\ell(A'(\overline K))$。）因此有必要更仔细地研究 $\pi$ 在 $U^{\mathrm{ef}}$ 上的\Term{action}{作用}。

按构造，这个\Term{action}{作用}在 $I$ 的\Term{normal-subgroup}{正规子群} $I'$ 上是平凡的，因而\Term{resolution}{分解}为\Term{quotient-group}{商群}
\[
 \pi'=\pi/I' \tag{4.1.3}
\]
\SourcePage{365}
的一个\Term{representation}{表示}。现在将对这一\Term{action}{作用}给出一个“与 $\ell$ 无关”的描述。
\textbf{4.2。}设 $K'$ 为 $K$ 的一个\Term{galois-extension}{Galois 扩张}，对应于 $\pi$ 的\Term{open-normal-subgroup}{开正规子群} $V$，且 $V\cap I=I'$。于是 $K'/K$ 的\Term{galois-group}{Galois 群} $G$ 可看作\Term{group-extension}{扩张}
\[
 1\longrightarrow\Gamma\longrightarrow G\longrightarrow G_\circ\longrightarrow1,
 \tag{4.2.1}
\]
其中
\[
 \Gamma=I/I' \tag{4.2.2}
\]
为\Term{inertia-subgroup}{惯性子群}，$G_\circ$ 为 $K'/K$ 的最大\Term{unramified-subextension}{非分歧子扩张} $K''$ 的\Term{galois-group}{Galois 群}。因此，\Term{exact-sequence}{正合列} (4.2.1) 是类似的\Term{exact-sequence}{正合列}
\[
 1\longrightarrow\Gamma\longrightarrow\pi'\longrightarrow\pi_\circ\longrightarrow1
 \tag{4.2.3}
\]
的像，其中 $\pi'$ 为 (4.1.3) 中\Term{action}{作用}于 $U^{\mathrm{ef}}$ 的\Term{group}{群}。于是有下述域的包含图：
\begin{equation}
\begin{tikzcd}[column sep=large,row sep=large]
 K \arrow[r,"G_\circ"] \arrow[dr,"G"'] \arrow[rr,bend left=26,"\pi_\circ"]
 & K''=\widetilde K\cap K' \arrow[r] \arrow[d,"\Gamma"]
 & \widetilde K \arrow[d,"\Gamma"] \arrow[dd,bend left=37,"I"] \\
 & K' \arrow[r] & \widetilde K' \arrow[d,"I'"] \\
 && \overline K
\end{tikzcd}\tag{4.2.4}
\end{equation}
设 $S'$ 为 $S$ 在 $K'$ 中的\Term{normalization}{正规化}，$A^*$ 为 $A_{K'}$ 的\Term{neron-model}{Néron 模型}。由\Term{semistable-reduction}{半稳定约化}的 Galois 判据（4.5），$A_{K'}$ 在 $S'$ 上有\Term{semistable-reduction}{半稳定约化}，即 $A_\circ^{*\circ}$ 是\Term{abelian-scheme}{阿贝尔概形}
\SourcePage{366}
被\Term{torus}{环面}的\Term{group-extension}{扩张}。另一方面，利用\Term{transport-of-structure}{结构搬运}，$G$ 以\Term{automorphism}{自同构}\Term{action}{作用}于 $A^*$，从而\Term{action}{作用}于 $A^{*\circ}$，并与 $G$ 在 $A^*$ 的\Term{base-scheme}{基概形} $S'$ 上的给定\Term{action}{作用}相容。于是显然，$\pi'$ 在 $U^{\mathrm{ef}}=T_\ell(A^{*\circ})^{\mathrm f}(\overline K)$ 上的\Term{action}{作用}，是由 $\pi'$ 在 $(A^{*\circ}/K',\overline K)$ 上的\Term{action}{作用}经\Term{transport-of-structure}{结构搬运}得到的；$\pi'$ 通过 $G$ 在 $A^*$ 上\Term{action}{作用}。考虑\Term{canonical-isomorphism}{典范同构}
\[
 U^{\mathrm{ef}}\xrightarrow{\sim}T_\ell(A^{*\circ})^{\mathrm f}(\overline K)
 \xrightarrow{\sim}T_\ell(A_\circ^{*\circ})(\overline k)
 \tag{4.2.6}
\]
（这对 $S'$ 上每个\Term{twisted-constant-l-adic-sheaf}{扭常值 $\ell$ 进层}都成立，参见第 I 讲\SGAsourcegap{原 PDF 第366页此处“cf. I”之后的交叉引注留空。}），便得到 $\pi'$ 在 $U^{\mathrm{ef}}$ 上\Term{action}{作用}的如下描述。

考虑 $S'$ 的\Term{residue-field}{剩余域} $k'\hookrightarrow\overline k$ 上的\Term{smooth}{光滑}\Term{connected}{连通}\Term{group-scheme}{群概形} $A_\circ^{*\circ}$；$k'$ 是 $k$ 的一个\Term{quasi-galois-extension}{拟 Galois 扩张}，其\Term{galois-group}{Galois 群}为 $G_\circ$。$\pi'$ 的\Term{quotient-group}{商群} $G$ \Term{action}{作用}于 $A_\circ^{*\circ}$，并与它通过 $G_\circ$ 在 $k'$ 上的\Term{action}{作用}相容。此外，$\pi'$ 通过其\Term{quotient}{商} $\pi_\circ$ \Term{action}{作用}于 $\overline k$，因而\Term{action}{作用}于 $(A_\circ^{*\circ}/k',\overline k)$ 这一情形。因此，利用\Term{transport-of-structure}{结构搬运}，$\pi'$ \Term{action}{作用}于
\[
 U^{\mathrm{ef}}=T_\ell(A_\circ^{*\circ})(\overline k),
\]
这就是所求的 $\pi'$ 的\Term{action}{作用}。特别地，$G$ 的\Term{subgroup}{子群} $\Gamma$ 在 $k'$ 上\Term{trivial-action}{平凡作用}，因而以 $k'$-\Term{automorphism}{自同构}\Term{action}{作用}于 $A_\circ^{*\circ}$；$\Gamma=I/I'$ 在 $T_\ell(A_\circ^{*\circ})(\overline k)$ 上的相应\Term{representation}{表示}，就是从\Term{inertia-group}{惯性群} $I$ 在 $U^{\mathrm{ef}}$ 上的自然\Term{representation}{表示}得到的\Term{representation}{表示}。

另一方面，$U$ 的\Term{essentially-toric-part}{本质环面部分} $U^{\mathrm{et}}$ 就对应于\Term{submodule}{子模}
\[
 U^{\mathrm{et}}\xrightarrow{\sim}T_\ell(T_\circ^*)(\overline k)
 \subset U^{\mathrm{ef}}\xrightarrow{\sim}T_\ell(A_\circ^{*\circ})(\overline k),
 \tag{4.2.7}
\]
\SourcePage{367}
其中 $T_\circ^*$ 是 $A_\circ^{*\circ}$ 的最大\Term{subtorus}{子环面}。同样，记
\[
 B_\circ^*=A_\circ^{*\circ}/T_\circ^* \tag{4.2.8}
\]
为 $A_\circ^{*\circ}$ 的\Term{abelian-part}{阿贝尔部分}。$G$ 在 $A_\circ^{*\circ}$ 上的\Term{action}{作用}分别诱导 $G$ 在 $T_\circ^*$ 与 $B_\circ^*$ 上的\Term{action}{作用}，而 $\pi'$ 在 $U^{\mathrm{et}}$ 与 $U^{\mathrm{ef}}/U^{\mathrm{et}}$ 上的\Term{action}{作用}，则借助\Term{isomorphism}{同构} (4.2.7) 及相应的\Term{isomorphism}{同构}
\[
 U^{\mathrm{ef}}/U^{\mathrm{et}}\xrightarrow{\sim}T_\ell(B_\circ^*)(\overline k)
 \tag{4.2.9}
\]
从前述\Term{action}{作用}得到。

由此及\Term{isogeny}{同源} (4.1.2)，对 $\pi'$（或 $\pi$）在 $U$ 上\Term{action}{作用}的\Term{character}{特征标} $g\mapsto\operatorname{Tr}g_U$，可得公式
\[
 \operatorname{Tr}g_{U_\ell}
 =\operatorname{Tr}g_{B_\circ^*,\ell}
 +\operatorname{Tr}g_{T_\circ^*,\ell}
 +\chi_\ell(g)\operatorname{Tr}g_{T_\circ^*,\ell}^{-1}.
 \tag{4.2.10}
\]
左边写成 $U_\ell$ 而不写 $U$，是为了提醒读者所选择的\Term{prime-number}{素数}为 $\ell\ne p$。这里 $g_{B_\circ^*,\ell}$ 与 $g_{T_\circ^*,\ell}$ 分别表示 $g$ 在\Term{tate-module}{Tate 模} $T_\ell(B_\circ^*)(\overline k)$ 与 $T_\ell(T_\circ^*)(\overline k)$ 上的\Term{action}{作用}，它们由刚才说明的 $g$ 在 $B_\circ^*$ 与 $T_\circ^*$ 上的\Term{action}{作用}得到；而
\[
 \chi_\ell:\pi'\longrightarrow\pi_\circ\longrightarrow\mathbb Z_\ell^*
 \tag{4.2.11}
\]
是由 $\mathbb Z_\ell(1)=T_\ell(\overline k^{\,*})$ 上的\Term{action}{作用}得到的 $\pi'$（或 $\pi_\circ$）的\Term{character}{特征标}，这个\Term{character}{特征标}在 $\Gamma^\circ$ 上是平凡的。

现在已可以证明 Serre--Tate 的一个猜想 [29，附录] 的一个特例\footnote{这实质上是所述关于 $H^i$ 的猜想的 $i=1$ 情形。}。
\SourcePage{368}
\begin{SGAresult}{\Term{theorem}{定理}}{4.3}
设 $S$ 为\Term{residue-field}{剩余域}为 $k$、\Term{function-field}{函数域}为 $K$ 的\Term{henselian}{Hensel}\Term{trait}{离散赋值环谱}，$A_K$ 为 $K$ 上的\Term{abelian-scheme}{阿贝尔概形}。考虑 $\pi=\operatorname{Gal}(\overline K/K)$ 在\Term{tate-module}{Tate 模} $U_\ell=T_\ell(A(\overline K))$ 上的\Term{representation}{表示}，其中 $\ell$ 是不同于\Term{residue-characteristic}{剩余特征}的\Term{prime-number}{素数}。

 a) 存在\Term{inertia-group}{惯性群}的一个\Term{open-subgroup}{开子群} $I'$，它在 $U_\ell$ 上作\Term{unipotent}{幺幂}\Term{action}{作用}，因此 $\pi$ 在 $U_\ell$ 上的\Term{representation}{表示}的\Term{character}{特征标}实际上是 $\pi/I'$ 上的函数。此外，这个函数在 $\Gamma=I/I'$ 上取整数值，且与 $\ell$ 无关。更\Term{generic}{一般}地，对每个 $g\in I$，$g_{U_\ell}$ 的\Term{characteristic-polynomial}{特征多项式}的系数都是与 $\ell$ 无关的有理整数。

 b) 假设 $k$ 为有 $q$ 个元素的\Term{finite-field}{有限域}，从而 $\pi_\circ=\pi/I\xrightarrow{\sim}\operatorname{Gal}(\overline k/k)$ 与 $\mathbb Z$ 的\Term{completion}{完备化} $\widehat{\mathbb Z}$ 等同，$\mathbb Z$ 的\Term{generator}{生成元} $1$ 对应于 $\overline k/k$ 的\Term{frobenius-automorphism}{Frobenius 自同构} $f_q:x\mapsto x^q$。则对每个在 $\pi_\circ$ 中的像属于 $\mathbb Z$ 的 $g\in\pi$，$g_{U_\ell}$ 的迹（更\Term{generic}{一般}地，其\Term{characteristic-polynomial}{特征多项式}的每个系数）都是与 $\ell$ 无关的有理数；若 $g$ 在 $\pi_\circ$ 中的像具有 $f_q^n$ 的\Term{form}{形式}，其中 $n\geq0$，则它甚至是整数。
\end{SGAresult}
\begin{proof}
例如来证明关于迹的断言。关于\Term{characteristic-polynomial}{特征多项式}系数的断言可用同样的方法证明；也可以指出，它们\Term{form}{形式}地由关于迹的断言\Term{pushout}{推出}：用各次迭代的迹计算这些系数（不过，这些公式引入了分母，因而不能给出所述的\Term{integrality}{整性}），再使用 Dwork--Fatou 的一个\Term{lemma}{引理}即可；
\SourcePage{369}
Kleiman 在 [14，2.8] 中给出了这个\Term{lemma}{引理}的一个漂亮证明。

 a) 使用公式 (4.2.10)，它在这里成为
\[
 \operatorname{Tr}g_{U_\ell}
 =\operatorname{Tr}g_{B_\circ^*,\ell}
 +\operatorname{Tr}g_{T_\circ^*,\ell}
 +\operatorname{Tr}g_{T_\circ^*,\ell}^{-1}. \tag{4.3.0}
\]
现在 $g_{B_\circ^*}$ 与 $g_{T_\circ^*}$ 分别为 $B_\circ^*$ 与 $T_\circ^*$ 的 $k'$-\Term{automorphism}{自同构}。因此只须证明相应的两个迹是与 $\ell$ 无关的有理整数；这是关于\Term{algebraic-group}{代数群}的\Term{generic}{一般}“几何”命题。对于 $g_{B_\circ^*}$，这是 Weil 关于\Term{abelian-variety}{阿贝尔簇}理论的一个著名结果 [5] [15]。对于 $g_{T_\circ^*}$，则几乎是显然的：借助\Term{dual-group}{对偶群} $M'$（\Term{isomorphic}{同构}于 $\mathbb Z^\mu$）上的相应\Term{representation}{表示}，将 $\Gamma$ 在 $T=T_\circ^*$ 上的\Term{representation}{表示}写出即可。由于有\Term{canonical-isomorphism}{典范同构}
\[
 T_\ell(T)\xrightarrow{\sim}\widetilde M_\ell(1)
 \mathrel{\mathop{=}^{\mathrm{d\text{é}f}}}
 M'^\vee\otimes_{\mathbb Z}\mathbb Z_\ell(1),
\]
所求的迹来自后一\Term{representation}{表示}的\Term{contragredient}{逆步表示}。

 b) 若 $g$ 在 $\pi_\circ$ 中的像为 $f_q^n$，则公式 (4.2.10) 变为
\[
 \operatorname{Tr}g_{U,\ell}
 =\operatorname{Tr}g_{B_\circ^*,\ell}
 +\operatorname{Tr}g_{T_\circ^*,\ell}
 +q^n\operatorname{Tr}g_{T_\circ^*,\ell}^{-1}. \tag{4.3.1}
\]
于是只须证明下述\Term{lemma}{引理}（其中令 $k\mapsto k'$，$A\mapsto A_\circ^{*\circ}\otimes_{k'}\overline k$）。
\end{proof}
\begin{SGAresult}{\Term{lemma}{引理}}{4.3.2}
设 $k$ 为一个\Term{finite-field}{有限域}的\Term{algebraic-closure}{代数闭包}，$A$ 为 $k$ 上\Term{finite-type}{有限型}的\Term{commutative-group-scheme}{交换群概形}，$g$ 为这个\Term{group-scheme}{群概形}的一个\Term{automorphism}{自同构}，它与 $\operatorname{Spec}(k)$ 的\Term{automorphism}{自同构} $f$ 相容；$f$ 由迭代的\Term{frobenius-automorphism}{Frobenius 自同构} $x\mapsto x^q$ 经\Term{transport-of-structure}{结构搬运}得到（其中 $q=p^r$ 是\Term{characteristic}{特征}的幂，$r\geq0$）。
\SourcePage{370}
因此，对 $g$ 在函数 $\varphi$ 上的\Term{action}{作用}，有
\[
 g^*(\lambda\varphi)=f^*(\lambda)g^*(\varphi)
 =\lambda^{q^{-1}}g^*(\varphi)
 =\lambda^{p^{-r}}g^*(\varphi)\qquad(\lambda\in k).
\]
考虑 $g$ 经“\Term{transport-of-structure}{结构搬运}”在 $T_\ell(A(k))$ 上诱导的\Term{automorphism}{自同构} $g_\ell$，其中 $\ell$ 是不等于 $p$ 的\Term{prime-number}{素数}。则 $g_\ell$ 的迹（更\Term{generic}{一般}地，其\Term{characteristic-polynomial}{特征多项式}的每个系数）是与 $\ell$ 无关的有理整数，其\Term{eigenvalue}{特征值}是绝对值等于 $q^{1/2}$ 或 $q$ 的\Term{algebraic-integer}{代数整数}。若 $A$ 为\Term{abelian-scheme}{阿贝尔概形}，则这些绝对值为 $q^{1/2}$；若 $A$ 为\Term{affine}{仿射}的，则\Term{eigenvalue}{特征值}具有 $qe_i$ 的\Term{form}{形式}，其中 $e_i$ 为\Term{root-of-unity}{单位根}。
\end{SGAresult}
\begin{proof}
考虑 $A$ 的底层绝对\Term{scheme}{概形}的“$q$ 次幂”\Term{endomorphism}{自同态} $f_q^A$，其定义为在点上取\Term{identity-map}{恒等映射}，在\Term{structure-sheaf}{结构层}上取 $q$ 次幂。这是\Term{group-scheme}{群概形} $A$ 的一个\Term{endomorphism}{自同态}，与基 $\operatorname{Spec}(k)$ 上的\Term{endomorphism}{自同态} $f_q^{\operatorname{Spec}(k)}=f^{-1}$ 相容。因此复合
\[
 F=f_q^A g \tag{4.3.3}
\]
是\Term{group-scheme}{群概形} $A$ 的一个与 $\operatorname{Spec}(k)$ 上\Term{identity-map}{恒等映射}相容的\Term{endomorphism}{自同态}，即\Term{group-scheme}{群概形} $A$ 的一个 $k$-\Term{endomorphism}{自同态}。我断言，由 $g$ 经\Term{transport-of-structure}{结构搬运}定义的 $A(k)$ 的\Term{automorphism}{自同构} $\Gamma_g$，即由公式
\[
 \Gamma_g(u)=g\circ u\circ f^{-1}\qquad\text{对 }u\in A(k), \tag{4.3.4}
\]
给出的\Term{automorphism}{自同构}，也可描述为
\[
 \Gamma_g(u)=F(u)\qquad\text{对 }u\in A(k). \tag{4.3.5}
\]
\SourcePage{371}
为此，只须看出，两边所表示的\Term{morphism}{态射} $\operatorname{Spec}(k)\to A$ 的底层\Term{set}{集合}映射相同；因此归结为证明如下等式：在 (4.3.5) 的左边左复合 $f_q^A$，在右边右复合 $f^{-1}$ 后得到的等式（因为 $f_q^A$ 与 $f^{-1}$ 在源\Term{scheme}{概形}的底层\Term{set}{集合}上都诱导\Term{identity-map}{恒等映射}）。而该公式正是 (4.3.3)、(4.3.4) 的直接推论。

将 (4.3.5) 应用于 $u\in{}_{\ell^\nu}A(k)$，再对 $\nu$ 取\Term{limit}{极限}，可得 $g_\ell$ 也\Term{solvable}{可解}释为与 $A$ 的\Term{endomorphism}{自同态} $F$ 相联系的 $T_\ell(A(k))$ 的\Term{endomorphism}{自同态} $T_\ell(F)$。另一方面，已知 $T_\ell(F)$ 的\Term{characteristic-polynomial}{特征多项式}具有与 $\ell$ 无关的整数系数：作\Term{devissage}{分拆}，化到 $A$ 为\Term{abelian-scheme}{阿贝尔概形}或\Term{torus}{环面}的情形，即得此结论；在这些情形中，该结果已经用于上面证明的 a) 部分。此外，结合\Term{abelian-scheme}{阿贝尔概形}情形下 Weil 的经典结果 [15]，这一证明也给出了关于\Term{eigenvalue}{特征值}的各项断言。

这就证明了 4.3.2，并完成了 4.3 的证明，同时也证明了以下推论。
\end{proof}
\begin{SGAresult}{推论}{4.4}
置身于 4.3 b) 的条件下。则对 $\pi'$ 中每个在 $\pi_\circ$ 中的像形如 $f_q^n$（$n\geq0$）的元素 $g$，在 $U^{\mathrm{et}}$、$U^{\mathrm{ef}}/U^{\mathrm{et}}$、$U/U^{\mathrm{ef}}$ 上所诱导的\Term{automorphism}{自同构}的\Term{eigenvalue}{特征值}都是\Term{algebraic-integer}{代数整数}，其绝对值分别为 $q^n$、$q^{n/2}$、$1$。对两端的部分，还可以更精确地说，相应的\Term{eigenvalue}{特征值}
\SourcePage{372}
分别为 $(q^ne_i)_{1\leq i\leq\mu}$ 和 $(e_i)_{1\leq i\leq\mu}$，其中 $e_i$（$1\leq i\leq\mu$）为\Term{root-of-unity}{单位根}。
\end{SGAresult}
\begin{proof}
事实上，对最后一个部分，公式 (4.1.2) 给出的\Term{eigenvalue}{特征值}为 $q^n/(q^ne_i)=e_i^{-1}$；但 $e_i$ 是\Term{root-of-unity}{单位根}，故 $e_i^{-1}=\overline e_i$（复\Term{conjugacy}{共轭}），且 $e_i$ 这一族在复\Term{conjugacy}{共轭}下\Term{stable}{稳定}（因为它是一个整数系数方程的解族）。所以 $e_i^{-1}$ 这一族也就是 $e_i$ 这一族，这就证明了 4.4 的最后一个断言。
\end{proof}
\textbf{4.5。}下面将 4.3 a) 应用于\Term{abelian-scheme}{阿贝尔概形}的\Term{local-conductor}{局部导子}理论。假设 $k$ \Term{algebraically-closed}{代数闭}\footnote{实际上只须 $k$ 为\Term{perfect-field}{完美域}；但换为 $S$ 的\Term{strict-henselization}{严格 Hensel 化}，便立即归结到所考虑的情形。}。回顾（SGA 5 X 6.2 或 [20]），对 $K$ 上每个\Term{etale}{平展}\Term{commutative-l-group-scheme}{交换 $\ell$ 群概形} $F_K$（$\ell\ne p$），按如下方式定义一个整数
\[
 \alpha_s(F_K)=\alpha_s(i_*(F_K)) \tag{4.5.1}
\]
（$i:\eta\to S$ 表示\Term{inclusion}{包含映射}）：选取 $K$ 的一个\Term{finite-galois-extension}{有限 Galois 扩张} $K'$，其\Term{group}{群}为 $\Gamma$，使 $F_K$ 在其上\Term{split}{分裂}，从而 $F_K$ 由一个有 $\Gamma$ \Term{action}{作用}的\Term{finite-l-group}{有限 $\ell$ 群} $M$ 定义。另一方面，对 $K$ 的\Term{extension}{扩张} $K'$，联系一个\Term{projective-module}{投射} $\mathbb Z_\ell[\Gamma]$-\Term{module}{模}（在相差\Term{isomorphism}{同构}的意义下定义）
\[
 \operatorname{Sw}=\operatorname{Sw}(K'/K,\mathbb Z_\ell), \tag{4.5.2}
\]
并置
\[
 \alpha_s(F_K)=\operatorname{long}\operatorname{Hom}_{\mathbb Z_\ell(\Gamma)}
 (\operatorname{Sw},M). \tag{4.5.3}
\]
可验证（同上），这一表达式不依赖于\Term{extension}{扩张} $K'/K$ 的选择。此外，由 $\operatorname{Sw}$ 为\Term{projective-module}{投射} $\mathbb Z_\ell[\Gamma]$-\Term{module}{模}，
\SourcePage{373}
可知 $\alpha_s(F_K)$ 对\Term{exact-sequence}{正合列}具有可\Term{additive}{加性}。

我们将把这一定义应用于\Term{sheaf}{层} ${}_\ell A_K$，其中 $A_K$ 是所考虑的 $K$ 上的\Term{abelian-scheme}{阿贝尔概形}。于是此时有整数
\[
 \alpha_s(A_K;\ell)=\operatorname{long}\operatorname{Hom}_{\mathbb Z_\ell(\Gamma)}
 (\operatorname{Sw},U_\ell\otimes_{\mathbb Z_\ell}\mathbb F_\ell),
 \qquad U_\ell=T_\ell(A(\overline K)), \tag{4.5.4}
\]
其中 $\Gamma$ 按上述方式选取（这一选择先验地依赖于 $\ell$）。整数 (4.5.4) 称为\Term{abelian-scheme}{阿贝尔概形} $A_K$ 关于\Term{discrete-valuation-ring}{离散赋值环} $V\subset K$ 的\Term{wild-conductor-exponent}{非驯导子指数}（或“野性\Term{conductor-exponent}{导子指数}”）\footnote{还应定义 $A_K$ 关于 $V$ 的\Term{tame-conductor-exponent}{驯导子指数}：它是 ${}_\ell A^\circ$ 的\Term{generic-geometric-fiber}{几何一般纤维}与\Term{special-geometric-fiber}{几何特殊纤维}在 $\mathbb Z/\ell\mathbb Z$ 上的\Term{length}{长度}之差。它与 $\ell$ 无关，用 (1.1.11) 的记号，等于 $2\lambda+\mu$。$A_K$ 的\Term{total-conductor-exponent}{（总）导子指数}定义为上述两个整数之和；它正是 SGA 5 X 6.2、[20] 中关于 ${}_\ell A^\circ$ 的 \Term{euler-poincare-formula}{Euler--Poincaré 公式}所含的\Term{local-term}{局部项}。}。当然，换为\Term{henselization}{Hensel 化}后，即使不假设 $V$ 为\Term{henselian}{Hensel}的，这个指数也有定义，而且可立即归结到 $S$ 为\Term{strictly-local}{严格局部}的情形。我们要证明：
\begin{SGAresult}{推论}{4.6}
\Term{wild-conductor-exponent}{非驯导子指数} (4.5.4) 不依赖于\Term{prime-number}{素数} $\ell\ne p$ 的选择。
\end{SGAresult}
\begin{proof}
可以假设 $S$ 为\Term{strictly-local}{严格局部}的。我们将计算 (4.5.4)，并将其中的 $\operatorname{Sw}$ 换成任意另一个\Term{projective-module}{投射} $\mathbb Z_\ell(\Gamma)$-\Term{module}{模} $L$。注意，$U_\ell$ 的\Term{filtration}{滤过} (4.1.1) 由显然是\Term{direct-summand}{直和因子}的 $\mathbb Z_\ell$-\Term{submodule}{子模}构成，故
\SourcePage{374}
这个\Term{filtration}{滤过}在 $U(\ell)=U_\ell\otimes_{\mathbb Z_\ell}\mathbb F_\ell$ 上诱导一个类似的\Term{filtration}{滤过}，其逐次因子由 (4.1.1) 中的因子 $V_i$（$i=1,2,3$）
\[
 U^{\mathrm{et}},\quad U^{\mathrm{ef}}/U^{\mathrm{et}},\quad U/U^{\mathrm{ef}}
\]
与 $\mathbb F_\ell$ 作\Term{tensor-product}{张量积}得到。当然，这个\Term{filtration}{滤过}在 $\Gamma$ 下\Term{stable}{稳定}，由可\Term{additive}{加性}给出 $\operatorname{long}\operatorname{Hom}_{\mathbb Z_\ell(\Gamma)}(L,U(\ell))$ 的三项\Term{resolution}{分解}，各项是以下\Term{group}{群}的\Term{length}{长度}：
\[
 \operatorname{Hom}_{\mathbb Z_\ell(\Gamma)}(L,V_i\otimes_{\mathbb Z_\ell}\mathbb F_\ell)
 \xrightarrow{\sim}
 \bigl(\operatorname{Hom}_{\mathbb Z_\ell(\Gamma)}(L,V_i)\bigr)
 \otimes_{\mathbb Z_\ell}\mathbb F_\ell. \tag{*}
\]
(*) 中写出的\Term{isomorphism}{同构}仍来自对 $L$ 的\Term{projective-module}{投射性}假设；这里只假设 $K'$ 选得\Term{large-relative-to-cardinal}{足够大}，使 $\Gamma$ \Term{action}{作用}于各 $V_i$。而 $\operatorname{Hom}_{\mathbb Z_\ell(\Gamma)}(L,V_i)$ 显然是\Term{torsion-free}{无挠} $\mathbb Z_\ell$-\Term{module}{模}，所以其\Term{rank}{秩}等于 (*) 右边的\Term{length}{长度}。设 $\varphi_i$ 为 $\Gamma$ 在 $V_i$ 中的\Term{representation}{表示}的\Term{character}{特征标}，$\psi$ 为 $\Gamma$ 在 $L$ 中的\Term{representation}{表示}的\Term{character}{特征标}，则得
\begin{align*}
 \operatorname{long}\operatorname{Hom}_{\mathbb Z_\ell(\Gamma)}(L,U(\ell))
 &=\sum_{1\leq i\leq3}N^{-1}\sum_{g\in\Gamma}\psi(g)\varphi_i(g)\\
 &=N^{-1}\sum_{g\in\Gamma}\psi(g)\varphi(g),
 \qquad N=\operatorname{card}\Gamma. \tag{4.6.1}
\end{align*}
这里 $\varphi(g)=\sum_i\varphi_i(g)$ 也就是从\Term{inertia-group}{惯性群} $I$ 上的函数，即 $I$ 在 $U$ 上的\Term{monodromy-representation}{单值性表示}的\Term{character}{特征标}，经取\Term{quotient}{商}得到的 $\Gamma$ 上的函数。由 4.3 a)，这个\Term{character}{特征标}是一个取有理整数值且与 $\ell$ 无关的函数。取 $L=\operatorname{Sw}$，$L$ 的\Term{character}{特征标} $\psi(g)$ 也如此，它等于 $a-N\delta+1$，其中 $\delta$ 是 $\Gamma$ 的“Dirac 函数”，$a$ 是 $K'/K$ 的“\Term{artin-conductor}{Artin 导子}”[26]。
\SourcePage{375}
这就证明了 4.6，同时还给出\Term{wild-conductor-exponent}{非驯导子指数}的以下与 $\ell$ 无关的表达式：
\[
 \alpha_s(A_K)=N^{-1}\sum_{g\in\Gamma}
 \bigl(a(g)-N\delta(g)+1\bigr)\varphi(g). \tag{4.6.2}
\]
其中 $\Gamma$ 为 $K$ 的一个\Term{finite-galois-extension}{有限 Galois 扩张} $K'$ 的\Term{galois-group}{Galois 群}，这个\Term{extension}{扩张}使 $U^{\mathrm{ef}}$ \Term{split}{分裂}（例如取这样的最小\Term{extension}{扩张}，参见 4.1）；$a(g)$ 为 $\Gamma$ 上的\Term{artin-conductor}{Artin 导子}，$\varphi(g)$ 为 (4.3.0) 中写出的 $\Gamma$ 上的迹函数。在\Term{semistable-reduction}{半稳定约化}情形中，可以取 $K'=K$，因而 $\Gamma=e$，得到
\[
 \alpha_s(A_K)=0\qquad\text{（半稳定约化情形）}. \tag{4.6.3}
\]
\end{proof}
最后，明确写出由 4.2 的\Term{generic}{一般}讨论得到的一个\Term{semistable-reduction}{半稳定约化}判据。
\begin{SGAresult}{命题}{4.7（M. Raynaud 的\Term{semistable-reduction}{半稳定约化}判据）}
设 $S$ 为\Term{trait}{离散赋值环谱}，$A_K$ 为其\Term{function-field}{函数域} $K$ 上定义的\Term{abelian-scheme}{阿贝尔概形}，$n>1$ 为与\Term{residue-characteristic}{剩余特征} $p$ 互素的整数\footnote{由 J. Tate 告诉我的一个例子可知，当 $n$ 等于\Term{residue-characteristic}{剩余特征} $p$ 时，4.7 不再成立，即便 $K$ 的\Term{characteristic}{特征}为零且 $A_K$ 是具有\Term{complex-multiplication}{复乘}的\Term{elliptic-curve}{椭圆曲线}，也不成立。}，并且 ${}_nA_K$ 关于 $S$ 是\Term{unramified}{非分歧}的，即\Term{inertia-group}{惯性群}在 ${}_nA(\overline K)$ 上\Term{trivial-action}{平凡作用}。则或者 $A_K$ 在 $S$ 上有\Term{semistable-reduction}{半稳定约化}，或者 $n=2$；在后一情形，若 $S$ 为\Term{strictly-local}{严格局部}的，则 $A_K$ 在 $K$ 的一个\Term{finite-galois-extension}{有限 Galois 扩张} $K'$ 上取得\Term{semistable-reduction}{半稳定约化}，其\Term{galois-group}{Galois 群}\Term{isomorphic}{同构}于 $(\mathbb Z/2\mathbb Z)^r$。
\end{SGAresult}
\SourcePage{376}
\begin{proof}
显然可以假设 $S$ 为\Term{strictly-local}{严格局部}的。用 4.1、4.2 的记号，取尽可能小的 $K'$，即假设 $\Gamma$ \Term{faithful-action}{忠实作用}于 $A_\circ^{*\circ}$。按假设，$\Gamma$ 在 ${}_nA(\overline K)$ 上\Term{trivial-action}{平凡作用}，因而在
\[
 {}_nA(\overline K)^{\mathrm{ef}}\xrightarrow{\sim}
 {}_nA_\circ^*(k)\supset{}_nA_\circ^{*\circ}(k)
\]
上也如此。于是由以下两个\Term{lemma}{引理}即可得出结论。
\end{proof}
\begin{SGAresult}{\Term{lemma}{引理}}{4.7.1（J.-P. Serre）}
设 $G$ 为域 $k$ 上的\Term{commutative-group-scheme}{交换群概形}，是\Term{abelian-scheme}{阿贝尔概形}被\Term{torus}{环面}的\Term{group-extension}{扩张}，$n\geq2$ 为整数。则 $G$ 的每个在 ${}_nG$ 上诱导\Term{identity-map}{恒等映射}的\Term{finite-order-automorphism}{有限阶自同构}，当 $n\ne2$ 时均为\Term{identity-map}{恒等映射}，当 $n=2$ 时为\Term{identity-map}{恒等映射}或二阶\Term{automorphism}{自同构}。
\end{SGAresult}
\begin{SGAresult}{\Term{lemma}{引理}}{4.7.2}
每个非\Term{identity-element}{单位元}素的阶均为 $2$ 的\Term{finite-group}{有限群}都\Term{isomorphic}{同构}于某个 $(\mathbb Z/2\mathbb Z)^r$，其中 $r\geq0$。
\end{SGAresult}
\textbf{4.7.3。}回顾 4.7.1 的众所周知的证明。注意，由于 $n\operatorname{id}_G$ 为\Term{epimorphism}{满态射}，$u$ 在 ${}_nG$ 上诱导\Term{identity-map}{恒等映射}，即 $1-u$ 在 ${}_nG$ 上诱导零映射，这一假设可表述为：在 $G$ 的\Term{endomorphism-ring}{自同态环}中有 $1-u=nv$。为证明所述结论，将 $u$ 生成的\Term{finite-cyclic-group}{有限循环群}\Term{resolution}{分解}为各个\Term{prime-number}{素数} $p$（遍历 $u$ 的阶的素因子）对应的 $p$-\Term{primary-component}{准素分量}，便化到 $u$ 的阶为\Term{prime-number}{素数}幂 $p^N$ 的情形。设 $E$ 为 $\operatorname{End}(G)$ 中由 $v$ 生成的\Term{subring}{子环}；它是 $\mathbb Z[T]/(T^{p^N}-1)$ 的\Term{quotient}{商}，因而在 $\mathbb Z$ 上\Term{finite-type}{有限型}，并且\Term{torsion-free}{无挠}（这\Term{form}{形式}地由 $G$ \Term{divisible}{可除}得出）。所以它\Term{embedding}{嵌入} $E_{\mathbb Q}=E\otimes_{\mathbb Z}\mathbb Q$，
\SourcePage{377}
后者是 $\mathbb Q[T]/(T^{p^N}-1)$ 的\Term{quotient}{商}；而该\Term{algebra}{代数}是 $\mathbb Q$ 的若干\Term{finite-extension}{有限扩张}域的\Term{direct-product}{直积}。故 $E_{\mathbb Q}$ 也如此。考虑 $E,u,v$ 在 $E_{\mathbb Q}$ 的各个域\Term{set-component}{分量} $K_i$ 中的像，便只须证明：若 $K=\mathbb Q(u)$ 是由一个 $p^m$ 次\Term{primitive-root-of-unity}{本原单位根} $u$ 生成的域，而整数 $n\geq2$ 满足 $u-1\equiv0\pmod n$（在 $K$ 的\Term{ring-of-integers}{整数环}中），则 $u=1$，或 $u=-1$ 且 $n=2$。这立即由如下众所周知的事实得到：若 $u\ne1$，即 $m\geq1$，则\Term{ideal}{理想} $(u-1)$ 为\Term{prime-ideal}{素理想}，且\Term{ideal}{理想} $(p)$ 是它的某个幂；除 $p=2,m=1$ 即 $u=-1$ 的情形外，幂指数均大于 $1$。事实上，$(u-1)$ 为\Term{prime-ideal}{素理想}且整除 $p$，所以 $n\geq2$ 能整除 $u-1$ 只可能在 $n=p$ 时发生；而 $p$ 整除 $u-1$ 只可能在 $(u-1)=(p)$ 即 $u=-1$ 时发生，故 $n=2$。

\textbf{4.7.4。}证明 4.7.2。只须证明 $\Gamma$ 交换。若 $x,y\in\Gamma$，则 $(xy)(xy)=1$；因 $x^2=1$、$y^2=1$，也可写成
\[
 xy=(xy)^{-1}=y^{-1}x^{-1}=yx,
\]
证毕。

\textbf{注记 4.7.5。}4.7.3 中的证明结合 4.7.2 表明：若 $G$ 为一个\Term{topos}{拓扑斯}中的\Term{divisible}{可除}\Term{commutative-group}{交换群}，$n\geq2$ 为整数，$\Gamma$ 为一个\Term{finite-group}{有限群}，它在 $G$ 上\Term{faithful-action}{忠实作用}，并在 ${}_nG$ 上诱导\Term{identity-map}{恒等映射}，则当 $n\geq3$ 时 $\Gamma$ 为\Term{unit}{单位}\Term{group}{群}，当 $n=2$ 时它\Term{isomorphic}{同构}于某个\Term{group}{群} $(\mathbb Z/2\mathbb Z)^r$。
